\subsection{Paramtric Surfaces and their Area}

\content{
}{% presentation
\vspace{4in}
}{% nopresentation
}{% comments
Define a parametric surface: 
$$ \vec{r}(u,v) = x(u,v)\vec{i}+y(u,v)\vec{j}+z(u,v)\vec{k}, $$
defined on a region in the $uv$-plane. 
}


\content{
\begin{example} Identify and sketch the surface whose vector equation is given
by 
$$ \vec{r}(u,v) = 2\cos u \vec{i} + v\vec{j} + 2\sin u\vec{k}. $$
\end{example}
}{% presentation
\vspace{3in}
}{% nopresentation
}{% comments
From here we will jump to the computer to see an example. 
}

\content{
\begin{example}
Find a vector function that represents the plane that passes through the point
$P_0$ with position vector $\vec{r}_0$ and that contains two nonparallel vectors
$\vec{a}$ and $\vec{b}$. 
\end{example}
}{% presentation
\vspace{3in}
}{% nopresentation
}{% comments
We will frequently need to take a given surface and find a parametrization for
it, rather than being given the parametrization. 
}

\content{
\begin{example} Find a prametric representation of the sphere 
$$ x^2+y^2+z^2 = a^2. $$
\end{example}
}{% presentation
\vspace{3in}
}{% nopresentation
}{% comments
}

\content{
\begin{example} Find a parametric representation of the cylinder 
$$ x^2+y^2 = 4,\quad 0\leq z\leq 1. $$
\end{example}
}{% presentation
\vspace{2in}
}{% nopresentation
}{% comments
Since this has a simple representation in cylindrical coordinates, we shall use
those. 
}

\subsubsection{Tangent Planes}

\content{

}{% presentation
\vspace{4in}
}{% nopresentation
}{% comments
Go over how the normal vector of the tagent plane is derived.  the normal vector
is given by 
$$ \vec{r}_u \times \vec{r}_v. $$
Mention that a surface is called smooth if this vector is never the zero vector. 
}

\content{
\begin{example} Find the tangent plane to the surface with parametric equations
$x=u^2$, $y=v^2$, and $z=u+2v$ at the point $(1,1,3)$.
\end{example}
}{% presentation
\vspace{3in}
}{% nopresentation
}{% comments
Normal vector is $\la -2u, -4u, 4uv\ra. $
}

\subsubsection{Surface Area}
\content{

}{% presentation
\vspace{4in}
}{% nopresentation
}{% comments
Go over the derivation of the Riemann Summ for the surface area
}

\content{
\begin{definition} If a smooth parametric surface $S$ is given by the equation 
$$ \vec{r}(u,v) = \la x(u,v) , y(u,v), z(u,v) \ra, $$
and $S$ is covered just once as $(u,v)$ ranges throughout a region $D$, then the
surface area of $S$ is
$$ A(S) =\iint_D |\vec{r}_u\times \vec{r}_v|\,dA. $$
\end{definition}
}{% presentation
}{% nopresentation
}{% comments
}

\content{
\begin{example} Find the surface area of a sphere of radius $a$.
\end{example}
}{% presentation
\vspace{4in}
}{% nopresentation
}{% comments
}

\content{
\begin{example} Find the surface area of the part of the paraboloid $z=x^2+y^2$
which lies under the plane $z=9$.
\end{example}
}{% presentation
\vspace{3in}
}{% nopresentation
}{% comments
Go over how for a regular surface parametrized by $x$ and $y$, we still use the
above fomula. (this is derived in the textbook as a separate formula to
remember)
}
